\section*{Capítulo \ref{ch:g5:stars-night-sky} --- Las estrellas y el cielo nocturno}

\begin{solution}{exo:g5:stars-night-sky:1}
Una \omterm{def:g4:solar-system:star}{estrella} es un sol: una bola gigantesca de materia
incandescente que fabrica su propia \omterm{def:g1:light-and-shadows:source}{luz}. Parecen chispas solo por
la distancia --- hasta la \omterm{def:g1:light-and-shadows:source}{luz} se pasa años en la carretera desde la
más cercana.
\end{solution}

\begin{solution}{exo:g5:stars-night-sky:2}
Sí --- brillan al mediodía exactamente igual que a medianoche. El
cielo diurno está inundado de \omterm{def:g1:light-and-shadows:source}{luz} del Sol dispersada por el \omterm{def:g2:air-around-us:air}{aire}, y
ninguna chispa puede lucir junto a ese resplandor, igual que ninguna
vela luce junto a una hoguera.
\end{solution}

\begin{solution}{exo:g5:stars-night-sky:3}
La blanco azulada --- los colores de las \omterm{def:g4:solar-system:star}{estrellas} van del naranja
(las más frías) al blanco azulado (las más calientes), pasando por el
amarillo. Una brasa que brilla naranja frente a un acero calentado al
blanco cuenta la misma historia.
\end{solution}

\begin{solution}{exo:g5:stars-night-sky:4}
Un dibujo con nombre formado por \omterm{def:g4:solar-system:star}{estrellas} brillantes, inventado como
mapa del cielo y ayuda para la memoria. Sus \omterm{def:g4:solar-system:star}{estrellas} no son vecinas
de verdad --- solo da la casualidad de que se alinean vistas desde la
Tierra.
\end{solution}

\begin{solution}{exo:g5:stars-night-sky:5}
Se busca el Carro; se toman las dos \omterm{def:g4:solar-system:star}{estrellas} punteras del borde
delantero de su cazo; se prolonga su línea unas cinco separaciones por
el cielo; la \omterm{def:g4:solar-system:star}{estrella} a la que se llega es la
\omterm{met:g5:stars-night-sky:northstar}{Estrella Polar}, y debajo de ella, en el horizonte, está el norte.
\end{solution}

\begin{solution}{exo:g5:stars-night-sky:6}
Está casi exactamente encima del eje de giro de la Tierra, así que el
giro del cielo la deja casi inmóvil --- una marca fija del norte que
sustituye a la \omterm{def:g4:magnets-and-compass:compass}{brújula}.
\end{solution}

\begin{solution}{exo:g5:stars-night-sky:7}
Lo que gira es la Tierra, que nos lleva consigo una vuelta al día. Ese
mismo giro hace desfilar al Sol por el cielo diurno --- un solo giro,
dos desfiles.
\end{solution}

\begin{solution}{exo:g5:stars-night-sky:8}
En incontables \omterm{def:g4:solar-system:star}{estrellas} tenues, demasiado apretadas para separarlas a
simple vista: la \omterm{ex:g5:stars-night-sky:milkyway}{Vía Láctea} --- nuestra propia ciudad plana de
\omterm{def:g4:solar-system:star}{estrellas}, vista de canto y desde dentro.
\end{solution}

\begin{solution}{exo:g5:stars-night-sky:9}
La Tierra que gira hace girar el cielo entero alrededor del eje que
pasa por la \omterm{met:g5:stars-night-sky:northstar}{Estrella Polar}. Cerca de ese punto quieto, las
\omterm{def:g4:solar-system:star}{estrellas} recorren círculos pequeños --- y la cámara registra
anillos. Hacia el este, el giro sube las \omterm{def:g4:solar-system:star}{estrellas} por encima del
horizonte siguiendo caminos inclinados --- y la cámara registra rayas
que suben. Un solo giro, dos maneras de verlo.
\end{solution}

\begin{solution}{exo:g5:stars-night-sky:10}
La dirección: las punteras del Carro le dan la \omterm{met:g5:stars-night-sky:northstar}{Estrella Polar} y,
con ella, el norte --- sin necesidad de aguja. El reloj: el propio
Carro gira alrededor de la \omterm{met:g5:stars-night-sky:northstar}{Estrella Polar} como la aguja de un reloj
enorme, una vuelta entera al día; mirar cuánto ha girado le dice cómo
va pasando la noche.
\end{solution}

\begin{solution}{exo:g5:stars-night-sky:11}
No: ves la \omterm{def:g4:solar-system:star}{estrella} tal como \emph{era}, hace años, cuando salió de
ella la \omterm{def:g1:light-and-shadows:source}{luz} de esta noche --- ves al mensajero, no al remitente.
Todos los telescopios miran, por lo tanto, hacia el pasado: cuanto más
lejos está la \omterm{def:g4:solar-system:star}{estrella}, más vieja es la noticia --- que es
precisamente lo que hace tan valiosa para los astrónomos la \omterm{def:g1:light-and-shadows:source}{luz}
lejana.
\end{solution}
